\scanpage{314}
\uline{Proposition.} Soit $X$ propre, plat, de prés. finie sur $S$. Alors $\uline{\mathrm{Pic}}^\tau_{X/S}$ est un sous-foncteur \uline{ouvert} \inserted{de prés. finie} de $\uline{\mathrm{Pic}}_{X/S}$. L'immersion ouverte $\uline{\mathrm{Pic}}^\tau_{X/S}\to\uline{\mathrm{Pic}}_{X/S}$ est aussi une immersion fermée si $X$ est \inserted{projectif} à fibres géom. intègres.
\margin{projectif est-il nécessaire ? \deleted{\unsure{plt. suffit}} \deleted{\uuline{\unsure{en fait}}}}

Cette assertion équivaut au

\uline{Corollaire.} $X$ comme dessus, $\uline{L}$ un module inversible. Alors l'ens \inserted{$U$} des $s\in S$ tels que $\uline{L}_s$ sur $X_s$ soit \deleted{\unsure{alg}} $\tau$-équiv. à $0$ est \uline{ouvert} \uline{rétrocompact}. Il est ouvert et fermé si $X$ est \inserted{projectif} à fibres géom. intègres.

On peut supposer $S$ noethérien. À prouver que $U$ est stable par générisation, et que si $s\in U$, alors $U\cap\overline{\{s\}}$ \deleted{est un \unsure{ouvert}} contient un ouvert non vide de $\overline{\{s\}}$. Or le critère (viii$'$ b$'$) implique facilement la première condition (en se ramenant au cas où \deleted{$U$ est \unsure{ouvert}} $S$ est un trait, et où son pt fermé $s$ est dans $U$). La deuxième condition est pratiquement triviale. Pour prouver que $U$ est \uline{fermé} si $X/S$ \inserted{projectif} à fibres géom. intègres, on est \deleted{\unsure{encore}} ramené à prouver qu'il est stable par spécialisation. On le voit sur le critère (viii e), \deleted{\unsure{a}}

\uline{Question} Éliminer l'hyp. projective dans la deuxième partie.
Prouver, si $X\to S$ est \uline{normal}, que $\uline{\mathrm{Pic}}^0_{X/S}$ est ``fermé'' dans $\uline{\mathrm{Pic}}_{X/S}$ mieux, que $\uline{\mathrm{Pic}}^0_{X/S}$ est \uline{propre sur $S$} ?? Cela signifie
\margin{?}
\figure{\begin{tikzpicture}\draw(0,0)--(0,1.2);\draw(.08,0)--(.08,1.2);\end{tikzpicture}}
\figure{\begin{tikzpicture}[x=1cm,y=1cm]
\draw[thin] (3.8,6.9)--(3.6,6.5);
\draw (3.3,5.95) .. controls (3.25,6.25) and (3.9,6.25) .. (4,6.35);
\draw (2.5,5.1) .. controls (2.48,5.35) and (2.9,5.34) .. (3.05,5.5);
\draw (.8,3.5) .. controls (.84,3.65) and (1.05,3.75) .. (1,3.9);
\draw (3.3,3.05) .. controls (3.27,3.3) and (3.75,3.3) .. (3.8,3.5);
\draw (.8,.5) .. controls (.8,.7) and (1.3,.7) .. (1.4,.9);
\end{tikzpicture}}
\mspage{313}
